\documentclass[12pt,a4paper]{ctexart}

% === 核心包 ===
\usepackage[backend=bibtex,style=numeric,sorting=none]{biblatex}
\usepackage{geometry}
\usepackage{graphicx}
\usepackage{booktabs}
\usepackage{longtable}
\usepackage{array}
\usepackage{calc}
\usepackage{hyperref}
\usepackage{xcolor}
\usepackage{enumitem}
\usepackage{etoolbox}
% longtable 用小号字，缓解 pandoc 表格列宽溢出
\AtBeginEnvironment{longtable}{\small}
% pandoc 输出的列表紧排命令（由 pandoc 生成但未定义）
\providecommand{\tightlist}{%
  \setlength{\itemsep}{0pt}\setlength{\parskip}{0pt}}
% 特殊 Unicode 符号映射（默认拉丁字体缺字）
\usepackage{newunicodechar}
\usepackage{amssymb}
% 星号与运算符
\newunicodechar{⭐}{\ensuremath{\bigstar}}
\newunicodechar{≤}{\ensuremath{\leq}}
\newunicodechar{≥}{\ensuremath{\geq}}
\newunicodechar{≈}{\ensuremath{\approx}}
% 希腊字母（热电主题正文常用）
\newunicodechar{β}{\ensuremath{\beta}}
\newunicodechar{κ}{\ensuremath{\kappa}}
\newunicodechar{σ}{\ensuremath{\sigma}}
% 下标
\newunicodechar{₀}{\textsubscript{0}}\newunicodechar{₁}{\textsubscript{1}}
\newunicodechar{₂}{\textsubscript{2}}\newunicodechar{₃}{\textsubscript{3}}
\newunicodechar{₄}{\textsubscript{4}}\newunicodechar{₈}{\textsubscript{8}}
\newunicodechar{₆}{\textsubscript{6}}\newunicodechar{₇}{\textsubscript{7}}
\newunicodechar{ₑ}{\textsubscript{e}}\newunicodechar{ₗ}{\textsubscript{l}}
\newunicodechar{ₓ}{\textsubscript{x}}
% 上标
\newunicodechar{⁺}{\ensuremath{^{+}}}
\newunicodechar{⁴}{\textsuperscript{4}}
\newunicodechar{₋}{\ensuremath{^{-}}}
\newunicodechar{ᵧ}{\textsuperscript{$\gamma$}}
% 带圈数字 ①–⑧
\newunicodechar{①}{\textcircled{\raisebox{-0.5pt}{1}}}
\newunicodechar{②}{\textcircled{\raisebox{-0.5pt}{2}}}
\newunicodechar{③}{\textcircled{\raisebox{-0.5pt}{3}}}
\newunicodechar{④}{\textcircled{\raisebox{-0.5pt}{4}}}
\newunicodechar{⑤}{\textcircled{\raisebox{-0.5pt}{5}}}
\newunicodechar{⑥}{\textcircled{\raisebox{-0.5pt}{6}}}
\newunicodechar{⑦}{\textcircled{\raisebox{-0.5pt}{7}}}
\newunicodechar{⑧}{\textcircled{\raisebox{-0.5pt}{8}}}
% 罗马数字
\newunicodechar{Ⅱ}{II}
% 状态图标
\newunicodechar{✅}{\checkmark}
\newunicodechar{⏳}{[TBD]}
\newunicodechar{⚠}{\textbf{/!\\}}
\geometry{margin=2.5cm}
\addbibresource{references.bib}

% === 元数据 ===
\title{ 文献调研报告 -- MOF 材料用于 CO2 捕获 }
\date{ 2026-08 }
\author{Pi-Agent（自主文献调研 Agent）}

\begin{document}
\maketitle
\sloppy  % 允许宽松换行，缓解长 URL/DOI 溢出

\section{文献调研报告 -- MOF 材料用于 CO2 捕获}\label{ux6587ux732eux8c03ux7814ux62a5ux544a-mof-ux6750ux6599ux7528ux4e8e-co2-ux6355ux83b7}

\textbf{调研主题}：金属有机框架（MOF）用于 CO2 捕集：材料-性质-构效关系系统综述 \textbf{生成日期}：2026-08-02 \textbf{文献规模}：546 篇（11 轮迭代，含历史记忆 152 篇 + Crossref/Sciverse 检索 394 篇），去重后约 370 篇有效 \textbf{数据文件}： - \texttt{gap\_report.md}：Research Gap 分析（10 项） - \texttt{knowledge\_graph.md}：知识图谱（33 条构效关系、4 个核心数值表） - \texttt{paper\_summaries.md}：论文摘要集（46 篇） - \texttt{audit\_trail.md}：API 调用审计证据链 - \texttt{discovery/discovery\_report.md}：阶段二构效关系发现报告 - \texttt{discovery/hypotheses.json}：假设定义与证据链

\begin{center}\rule{0.5\linewidth}{0.5pt}\end{center}

\subsection{摘要}\label{ux6458ux8981}

本报告系统梳理了 MOF 用于 CO2 捕获的文献前沿，覆盖六个核心材料体系（MOF-74 系列、ZIF 系列、MIL 系列、UiO/HKUST 系列、dobpdc 胺功能化系列、MOF 复合/载体体系）与四大性质维度（吸附容量、选择性、等量吸附热 Qst、水/杂质稳定性）。经 11 轮自主迭代（首轮 152 篇历史记忆 + 10 轮增量检索 394 篇），Agent 自主识别 10 项 Research Gap（Gap 1-10），构建 33 条关键构效关系（R1-R33），生成 5 条定量假设，完成贝叶斯/MCTS 搜索验证与 Materials Project/OQMD 外部交叉验证。

最突出的科学问题在于： 1. \textbf{双金属 MOF 比例-容量定量关系缺失}（Gap 1）：倒 U 型规律存在但非普适，受金属组合电子结构控制 2. \textbf{水-CO2 竞争/协同机理矛盾}（Gap 2）：经多轮迭代统一为 OMS 型（抑制）vs 胺型（促进）二分机理 3. \textbf{缺陷类型对 OMS/容量的方向性控制}（Gap 9）：缺失配体型提升容量，溶剂甲酸盐占据型降低容量

这些 Gap 已映射到阶段二的 5 条构效关系假设，其中 4 条置信度 \textgreater= 0.7 且具备独立证据链支撑。

\begin{center}\rule{0.5\linewidth}{0.5pt}\end{center}

\subsection{一、引言与研究背景}\label{ux4e00ux5f15ux8a00ux4e0eux7814ux7a76ux80ccux666f}

\subsubsection{1.1 CO2 捕获的材料需求}\label{co2-ux6355ux83b7ux7684ux6750ux6599ux9700ux6c42}

CO2 捕集、利用与封存（CCUS）是实现碳中和目标的核心技术环节。在捕集环节中，固体吸附剂相比传统胺溶液具有再生能耗低、无腐蚀、操作灵活等优势。金属有机框架（MOF），凭借其高比表面积（可达 5000+ m2/g）、孔径可调（微孔至介孔）、金属节点与有机配体的组合空间近乎无穷，成为吸附法捕集中最有前景的材料类别。

\subsubsection{1.2 MOF 用于 CO2 捕获的核心科学问题}\label{mof-ux7528ux4e8e-co2-ux6355ux83b7ux7684ux6838ux5fc3ux79d1ux5b66ux95eeux9898}

本批文献（546 篇，11 轮检索）揭示了三个结构化的科学问题：

{\def\LTcaptype{none} % do not increment counter
\begin{longtable}[]{@{}
  >{\raggedright\arraybackslash}p{(\linewidth - 4\tabcolsep) * \real{0.2143}}
  >{\raggedright\arraybackslash}p{(\linewidth - 4\tabcolsep) * \real{0.3214}}
  >{\raggedright\arraybackslash}p{(\linewidth - 4\tabcolsep) * \real{0.4643}}@{}}
\toprule\noalign{}
\begin{minipage}[b]{\linewidth}\raggedright
问题
\end{minipage} & \begin{minipage}[b]{\linewidth}\raggedright
具体表现
\end{minipage} & \begin{minipage}[b]{\linewidth}\raggedright
文献证据示例
\end{minipage} \\
\midrule\noalign{}
\endhead
\bottomrule\noalign{}
\endlastfoot
\textbf{知识碎片化} & 吸附容量、选择性、Qst 等关键性能分散在上百篇论文中，缺乏统一结构化整理 & 546 篇文献的性能数据需 Agent 自行组织为知识图谱 \\
\textbf{结论矛盾} & 同一材料在不同合成条件下性能数值差异巨大；水对 CO2 吸附的影响方向存在矛盾 & Ni-MOF-74 的 CO2 容量报道值横跨 3.99--8.29 mmol/g（\cite{p20} vs \cite{p62}）；水促进（\cite{p16}）vs 无影响（\cite{p129}）vs 无竞争（\cite{p139}）vs 衰减（\cite{p128}） \\
\textbf{权衡关系未被刻画} & 高容量、高选择性、低再生能耗三者存在公认 trade-off，但 Pareto 前沿形状与驱动因素无定量结论 & Qst 报道值 25--40 kJ/mol 窗口 vs 29 kJ/mol 甜点（\cite{p27}）；Gap 3 定量 \\
\end{longtable}
}

\subsubsection{1.3 本调研的范围与结构}\label{ux672cux8c03ux7814ux7684ux8303ux56f4ux4e0eux7ed3ux6784}

本调研分为四个层次：（1）材料体系分类与关键性能数值整理；（2）构效关系（R1-R33）的识别与组织；（3）Research Gap 的发现、证据收集与优先级排序；（4）对阶段二（构效关系假设生成与搜索验证）的指引。

\begin{center}\rule{0.5\linewidth}{0.5pt}\end{center}

\subsection{二、检索策略与方法}\label{ux4e8cux68c0ux7d22ux7b56ux7565ux4e0eux65b9ux6cd5}

\subsubsection{2.1 检索流程}\label{ux68c0ux7d22ux6d41ux7a0b}

Agent 采用 11 轮迭代的自主检索策略：

{\def\LTcaptype{none} % do not increment counter
\begin{longtable}[]{@{}
  >{\raggedright\arraybackslash}p{(\linewidth - 6\tabcolsep) * \real{0.1818}}
  >{\raggedright\arraybackslash}p{(\linewidth - 6\tabcolsep) * \real{0.2424}}
  >{\raggedright\arraybackslash}p{(\linewidth - 6\tabcolsep) * \real{0.3030}}
  >{\raggedright\arraybackslash}p{(\linewidth - 6\tabcolsep) * \real{0.2727}}@{}}
\toprule\noalign{}
\begin{minipage}[b]{\linewidth}\raggedright
轮次
\end{minipage} & \begin{minipage}[b]{\linewidth}\raggedright
数据源
\end{minipage} & \begin{minipage}[b]{\linewidth}\raggedright
新增文献数
\end{minipage} & \begin{minipage}[b]{\linewidth}\raggedright
聚焦方向
\end{minipage} \\
\midrule\noalign{}
\endhead
\bottomrule\noalign{}
\endlastfoot
首轮 & 历史记忆（Sciverse + arXiv） & 152 & MOF CO2 捕获全景 \\
第二轮 & Crossref API（20 组检索词） & 224 & 核心捕获、OMS、胺功能化、双金属、水稳定、DAC、TSA/VPSA、Qst、ML、MIL/ZIF/UiO \\
第三轮 & Sciverse 专项 & 45 & 杂质气体与湿气抗性、ML 筛选 \\
第四轮 & Sciverse 双金属专题（4 组检索词） & 84 & 双金属 MOF-74 混合金属比例 / 掺杂比例 OMS / DFT 电子结构 / CO2-N2 选择性 Qst \\
第五轮 & Sciverse s-d 组合专题（2 组检索词） & 10 & MgCu/MgNi/MgZn 体系 + 轨道类型假说检验 \\
第六轮 & Sciverse + Crossref （3 组检索词） & 26 & 组成可控性（NMR 配分解析 + 机械化学合成）+ DFT 方法学 \\
第七--十一轮 & 增量检索，策略持续优化 & \textasciitilde20 & dobpdc 胺变体体系 + 缺陷-水交叉旁证 + 辫状链机理 \\
\end{longtable}
}

\textbf{总计}：546 篇文献，覆盖 8 条检索主线，检索词 40+ 组。

\subsubsection{2.2 数据源与审计}\label{ux6570ux636eux6e90ux4e0eux5ba1ux8ba1}

{\def\LTcaptype{none} % do not increment counter
\begin{longtable}[]{@{}
  >{\raggedright\arraybackslash}p{(\linewidth - 4\tabcolsep) * \real{0.2759}}
  >{\raggedright\arraybackslash}p{(\linewidth - 4\tabcolsep) * \real{0.3103}}
  >{\raggedright\arraybackslash}p{(\linewidth - 4\tabcolsep) * \real{0.4138}}@{}}
\toprule\noalign{}
\begin{minipage}[b]{\linewidth}\raggedright
数据源
\end{minipage} & \begin{minipage}[b]{\linewidth}\raggedright
接入方式
\end{minipage} & \begin{minipage}[b]{\linewidth}\raggedright
审计记录位置
\end{minipage} \\
\midrule\noalign{}
\endhead
\bottomrule\noalign{}
\endlastfoot
arXiv & REST API（免费） & \texttt{workspace/data/literature\_cache/search\_log.jsonl} \\
Crossref & REST API（免费） & 同上 \\
Sciverse & MCP / Skill / REST 三级自适应 & \texttt{workspace/logs/sciverse\_skill\_log.jsonl} \\
\end{longtable}
}

所有 API 调用均生成标准化审计记录（详见 \texttt{audit\_trail.md}），包括调用 ID、参数哈希、时间戳、结果摘要，支持从任意 Gap/假设中的论文引用（p\# 编号）逐层追溯到原始 API 请求。

\subsubsection{2.3 PDF 解析}\label{pdf-ux89e3ux6790}

采用双引擎策略： - \textbf{MinerU Cloud v1}：远程 API，对学术论文 PDF 的表格/公式/参考文献解析精度高 - \textbf{markitdown}（本地回退引擎）：当 MinerU 不可用时，本地解析 HTML/text 格式

解析结果汇总于 \texttt{paper\_summaries.md}（46 篇完整摘要，含标题、DOI、摘要片段）。

\begin{center}\rule{0.5\linewidth}{0.5pt}\end{center}

\subsection{三、知识图谱总览}\label{ux4e09ux77e5ux8bc6ux56feux8c31ux603bux89c8}

\begin{quote}
完整图谱见 \texttt{knowledge\_graph.md}。本节提取核心内容供快速阅读。
\end{quote}

\subsubsection{3.1 材料体系分类（6 大类，\textasciitilde55 种）}\label{ux6750ux6599ux4f53ux7cfbux5206ux7c7b6-ux5927ux7c7b55-ux79cd}

\paragraph{MOF-74 系列（OMS 型）-- 证据最密集}\label{mof-74-ux7cfbux5217oms-ux578b-ux8bc1ux636eux6700ux5bc6ux96c6}

MOF-74 系是 CO2 捕获领域研究最密集（\textasciitilde120 篇直接文献），因其一维六角孔道中的开放金属位点（OMS）提供了最高的低分压 CO2 容量。

{\def\LTcaptype{none} % do not increment counter
\begin{longtable}[]{@{}
  >{\raggedright\arraybackslash}p{(\linewidth - 6\tabcolsep) * \real{0.2500}}
  >{\raggedright\arraybackslash}p{(\linewidth - 6\tabcolsep) * \real{0.2500}}
  >{\raggedright\arraybackslash}p{(\linewidth - 6\tabcolsep) * \real{0.2500}}
  >{\raggedright\arraybackslash}p{(\linewidth - 6\tabcolsep) * \real{0.2500}}@{}}
\toprule\noalign{}
\begin{minipage}[b]{\linewidth}\raggedright
材料
\end{minipage} & \begin{minipage}[b]{\linewidth}\raggedright
CO2 容量 (mmol/g)
\end{minipage} & \begin{minipage}[b]{\linewidth}\raggedright
Qst (kJ/mol)
\end{minipage} & \begin{minipage}[b]{\linewidth}\raggedright
备注
\end{minipage} \\
\midrule\noalign{}
\endhead
\bottomrule\noalign{}
\endlastfoot
Mg-MOF-74 & 3.67--8.6 & 42--47 & 容量最高但水敏（R7）; 晶间空间影响吸附（\cite{p6} GCMC） \\
Ni-MOF-74 & \textasciitilde3.99--8.29 & 可调 & 容量矛盾区间大（合成/后处理敏感，R26） \\
NiCo-MOF-74 & 8.30 & -- & 双金属协同 \textgreater{} 单金属（\cite{p62}, R5） \\
Co-MOF-74 & 5.03 & -- & 微波合成（\cite{p189}） \\
Fe-MOF-74 & -- & 36 & Mott 绝缘体挑战 DFT 泛函（六轮） \\
缺陷 MOF-74 & 缺陷提升 & -- & HBDC 缺陷产生额外 OMS（\cite{p7}/\cite{p13}/\cite{p192}, R11） \\
M-MOF-74 多金属 & -- & -- & Mg/Mn/Fe/Ni/Zn CO2 结合势能景观（\cite{p17}/\cite{p185}） \\
\end{longtable}
}

\textbf{双金属 MOF-74 体系}：CoMn（\cite{p65}, \textbf{最佳 1:1}）、Fe/Cu（\cite{p67}, 双金属 \textgreater{} 单金属）、Ni-Cu-MOF-74 MMM（\cite{p22}, 3188 Barrer / 35.1 选择性）、机械化学法 12 种 1:1 双金属（R23）。

\paragraph{胺功能化 dobpdc 系列 -- DAC 前沿}\label{ux80faux529fux80fdux5316-dobpdc-ux7cfbux5217-dac-ux524dux6cbf}

{\def\LTcaptype{none} % do not increment counter
\begin{longtable}[]{@{}
  >{\raggedright\arraybackslash}p{(\linewidth - 4\tabcolsep) * \real{0.2857}}
  >{\raggedright\arraybackslash}p{(\linewidth - 4\tabcolsep) * \real{0.4286}}
  >{\raggedright\arraybackslash}p{(\linewidth - 4\tabcolsep) * \real{0.2857}}@{}}
\toprule\noalign{}
\begin{minipage}[b]{\linewidth}\raggedright
材料
\end{minipage} & \begin{minipage}[b]{\linewidth}\raggedright
关键特征
\end{minipage} & \begin{minipage}[b]{\linewidth}\raggedright
文献
\end{minipage} \\
\midrule\noalign{}
\endhead
\bottomrule\noalign{}
\endlastfoot
mmen-Mg2(dobpdc) & 台阶等温线，2.0 mmol/g @0.39 mbar/25 C（经典 DAC） & \cite{r11s1_a7d3df468203} \\
dmpn-Mg2(dobpdc) & 煤烟气混合机理，2.42 mmol/g 工作容量 & \cite{r11s0_9932dacb5791} \\
双胺接枝 MOF（diamine-appended） & 低 RH 下 CO2 吸附增强，协同晶格理论 & \cite{r10s0_7e6841cd666f}（R29） \\
H2O/CO2 共吸附''辫状链''（braided chain） & ab initio + 统计力学预测新构型 & \cite{r11s2_1fdf67d949bc}（R32） \\
IRMOF-74-III-CH2NH2 & 65\% RH 湿态选择性捕获（OMS 型外的湿态实证） & \cite{r11s4_2c91341a5f47} \\
\end{longtable}
}

\paragraph{其他重要体系}\label{ux5176ux4ed6ux91cdux8981ux4f53ux7cfb}

\begin{itemize}
\tightlist
\item
  \textbf{ZIF 系列}：ZIF-8（\cite{p30} GCMC）、ZIF-8@Zn-MOF-74 核壳（\cite{p2}）、CoNi-ZIF-67 MMM（\cite{p75}）
\item
  \textbf{MIL 系列}：MIL-101(Cr)（\cite{p28}）、MIL-101(Cr,Mg) 双金属（\cite{p147}, R5）、MIL-100(Fe)（\cite{p166} IAST vs 共吸附）、氨基-MIL-53（\cite{p156}）
\item
  \textbf{UiO / HKUST 系列}：UiO-66（\cite{p30}, 深共晶溶剂接枝 \cite{p25}）、Cu-BTC/HKUST-1（\cite{p28}/\cite{p173}）、胺功能化 UiO-66 水桥替代缺陷（R30）
\item
  \textbf{特殊 MOF}：MOF-177（AAIL 封装 3x 容量, \cite{p104}）、MOF-303（机械化学 9.5 mmol/g@0 C, \cite{p76}）、ELM-11（NO2 杂质影响 gate-opening, \cite{p216}）
\item
  \textbf{载体/对比体系}：SBA-15/MCM-41/KIT-6 胺接枝、沸石 13X、硅胶 PEI、多孔碳
\end{itemize}

\subsubsection{3.2 性质维度（10 项）}\label{ux6027ux8d28ux7ef4ux5ea610-ux9879}

\begin{enumerate}
\def\labelenumi{\arabic{enumi}.}
\tightlist
\item
  \textbf{CO2 吸附容量}（mmol/g）：Mg-MOF-74 3.67--8.6；MOF-303 9.5@0 C；胺改性 Mg-MOF-74 3.9 饱和/1.27 动态
\item
  \textbf{CO2/N2 选择性}：Ni-Cu-MOF-74 MMM 35.1；MOF-177 AAIL 5-\textgreater13
\item
  \textbf{等量吸附热 Qst}（kJ/mol）：Mg-MOF-74 42--47；MOF-74(Ni) 可调（\cite{p27}）；SBA-15\_APTES 胺化（\cite{p129}）
\item
  \textbf{水稳定性/湿气行为}：Mg-MOF-74 水敏；核壳调控（\cite{p210}）；疏水化（\cite{p201}/\cite{p224}）；胺型 RH 峰值增强（R29）
\item
  \textbf{循环稳定性/再生性}：Mg-MOF-74 10 循环稳定（\cite{p197}）；PEI 硅胶 20 循环衰减（\cite{p206}）
\item
  \textbf{BET 比表面积}：Mg-MOF-74 提升（\cite{p197}）；MOF-303 1180 m2/g（\cite{p76}）
\item
  \textbf{超微孔体积}：胺改性 Mg-MOF-74 0.46 cm3/g（\cite{p193}）
\item
  \textbf{OMS 密度/缺陷密度}：缺陷 MOF-74 额外 OMS（\cite{p7}/\cite{p13}/\cite{p192}, R1/R11/R25）
\item
  \textbf{再生能耗/工艺能效}：TSA/VPSA Mg-MOF-74 CFD（\cite{p176}/\cite{p178}, R12）
\item
  \textbf{杂质气体耐受性}：NO2 对 ELM-11（\cite{p216}, Gap 8）
\end{enumerate}

\subsubsection{3.3 关键构效关系（R1-R33，节选）}\label{ux5173ux952eux6784ux6548ux5173ux7cfbr1-r33ux8282ux9009}

{\def\LTcaptype{none} % do not increment counter
\begin{longtable}[]{@{}
  >{\raggedright\arraybackslash}p{(\linewidth - 8\tabcolsep) * \real{0.1667}}
  >{\raggedright\arraybackslash}p{(\linewidth - 8\tabcolsep) * \real{0.1667}}
  >{\raggedright\arraybackslash}p{(\linewidth - 8\tabcolsep) * \real{0.1667}}
  >{\raggedright\arraybackslash}p{(\linewidth - 8\tabcolsep) * \real{0.2500}}
  >{\raggedright\arraybackslash}p{(\linewidth - 8\tabcolsep) * \real{0.2500}}@{}}
\toprule\noalign{}
\begin{minipage}[b]{\linewidth}\raggedright
编号
\end{minipage} & \begin{minipage}[b]{\linewidth}\raggedright
关系
\end{minipage} & \begin{minipage}[b]{\linewidth}\raggedright
方向
\end{minipage} & \begin{minipage}[b]{\linewidth}\raggedright
证据强度
\end{minipage} & \begin{minipage}[b]{\linewidth}\raggedright
关键文献
\end{minipage} \\
\midrule\noalign{}
\endhead
\bottomrule\noalign{}
\endlastfoot
R1 & OMS 密度 -\textgreater{} CO2 容量 & 正相关 & 高 & \cite{p10}/\cite{p13}/\cite{p192} \\
R3 & Qst vs 再生能耗 & 权衡 & 高 & \cite{p27}/\cite{p176}/\cite{p178} \\
R5 & 双金属比例 -\textgreater{} 容量 & \textbf{倒 U 型（1:1 最优），非普适} & 高 & \cite{p65}/\cite{p67}/\cite{p147} \\
R7 & 水蒸气 -\textgreater{} OMS 型容量 & 负相关 & 高 & \cite{p186}/\cite{p182}/\cite{p196} \\
R8 & 水蒸气 -\textgreater{} 胺型容量 & 先升后降（峰值 30--50\% RH） & 中-\textgreater 高 & \cite{p224}/\cite{p210}/\cite{p201} \\
R11 & 缺陷工程 -\textgreater{} OMS -\textgreater{} 容量 & 正相关（缺失配体型） & 高 & \cite{p7}/\cite{p13}/\cite{p192} \\
R19 & 双金属曲线形状 -\textgreater{} 金属电子结构（d-d 倒U vs s-d 单调） & 类型依赖 & 中 & scv:ecfe34f94197 \\
R22 & 双金属配分模式 -\textgreater{} 非随机 & NMR 实证 & 高 & 六轮 \\
R25 & 缺陷类型 -\textgreater{} 容量方向 \textbf{二分} & 缺失配体型 + vs 溶剂甲酸盐型 - & 高 & Nature Commun 2023 \\
R29 & 胺功能化 + 低 RH -\textgreater{} CO2 容量 & 增强（协同晶格） & 高 & \cite{r10s0_7e6841cd666f} + 十一轮辫状链 \\
R30 & 水桥 -\textgreater{} 缺陷替代 -\textgreater{} 容量 & 增强（Gap 2 x Gap 9 交叉） & 高 & \cite{r10s0_1fbbe940f717} + 十一轮旁证 \\
R32 & H2O+CO2 共吸附 -\textgreater{} 辫状链构型 & 原子尺度机理 & 中-高 & \cite{r11s2_1fdf67d949bc} \\
\end{longtable}
}

完整 R1-R33 详见 \texttt{knowledge\_graph.md}。

\subsubsection{3.4 核心数值表}\label{ux6838ux5fc3ux6570ux503cux8868}

{\def\LTcaptype{none} % do not increment counter
\begin{longtable}[]{@{}llllll@{}}
\toprule\noalign{}
材料 & CO2 容量 (mmol/g) & Qst (kJ/mol) & 选择性 & 条件 & 证据 \\
\midrule\noalign{}
\endhead
\bottomrule\noalign{}
\endlastfoot
Mg-MOF-74 & 8.6 & 42--47 & \textasciitilde180 & 298K/1bar & \cite{p20} \\
Mg-MOF-74（动态） & 3.67 & -- & -- & 30 C/0.1bar & \cite{p197} \\
胺-Mg-MOF-74 & 3.9（饱和）/1.27（动态） & -- & -- & 25 C & \cite{p193} \\
MOF-303 & 9.5 & -- & -- & 0 C/1bar & \cite{p76} \\
MOF-74(Ni) & \textasciitilde8.29 & 可调 & -- & 298K/1bar & \cite{p20} \\
NiCo-MOF-74 & 8.30 & -- & -- & 298K/1bar & \cite{p62} \\
Ni-MOF-74（优化洗涤） & 5.80 & -- & -- & 1bar/25 C & R26 \\
Ni-MOF-74（基础合成） & 3.99 & -- & -- & 298K/1bar & \cite{p62} \\
Co-MOF-74 & 5.03 & -- & -- & 298K/1bar & \cite{p62} \\
Ni-Cu-MOF-74 MMM & -- & -- & 35.1（3188 Barrer） & 膜 & \cite{p22} \\
AAIL@MOF-177 & 3x 原始 & -- & 5-\textgreater13 & 0.2bar/303K & \cite{p104} \\
ELM-11（NO2 暴露后） & 减少 & -- & -- & -- & \cite{p216} \\
\end{longtable}
}

\begin{center}\rule{0.5\linewidth}{0.5pt}\end{center}

\subsection{四、Research Gap 分析}\label{ux56dbresearch-gap-ux5206ux6790}

\begin{quote}
完整分析见 \texttt{gap\_report.md}（10 项 Gap，含优先级排序）。本节按严重程度从高到低概述。
\end{quote}

\subsubsection{4.1 高严重程度 Gap（Top 3）}\label{ux9ad8ux4e25ux91cdux7a0bux5ea6-gaptop-3}

\paragraph{Gap 1：双金属 MOF 金属比例-容量定量关系缺失}\label{gap-1ux53ccux91d1ux5c5e-mof-ux91d1ux5c5eux6bd4ux4f8b-ux5bb9ux91cfux5b9aux91cfux5173ux7cfbux7f3aux5931}

\begin{itemize}
\tightlist
\item
  \textbf{类型}：缺失连接（定量关系），\textbf{置信度}：0.95
\item
  \textbf{核心发现}：CoMn-MOF-74 最佳 Co/Mn = \textbf{1:1}（\cite{p65}）；Fe/Cu-MOF 双金属 \textgreater{} 单金属（\cite{p67}）；Ni-Cu-MOF-74 MMM (\cite{p22})。但 Cu/Mg-MOF-74（s-d 组合）吸附单调增，非倒 U（R19）。倒 U 型适用于 d-d 组合但不普适。
\item
  \textbf{合成敏感性矛盾}：Ni-MOF-74 容量 3.99（\cite{p62}）vs 8.29（\cite{p20}）vs 5.80（R26, 洗涤优化）--合成/后处理主导有效容量（R26）。
\item
  \textbf{验证方案}：梯度金属比合成 + GCMC/DFT + ICP 实测组成 + 后处理标准化
\end{itemize}

\paragraph{Gap 2：水-CO2 竞争/协同机理矛盾}\label{gap-2ux6c34-co2-ux7adeux4e89ux534fux540cux673aux7406ux77dbux76fe}

\begin{itemize}
\tightlist
\item
  \textbf{类型}：矛盾结论--机理统一，\textbf{置信度}：0.97
\item
  \textbf{核心发现}：经 11 轮迭代，矛盾统一为 \textbf{OMS 型（水占据 OMS--抑制）vs 胺型（水参与链形成--促进）} 二分。胺分支获原子尺度支撑：辫状链（braided chain）构型（R32, \cite{r11s2_1fdf67d949bc}）+ 协同晶格理论（R29）+ diamine 多体系实证（十一轮闭环）。
\item
  \textbf{验证方案}：统一 0--90\% RH breakthrough 矩阵（OMS 型 vs 胺型对照），量化胺型峰值 RH 与 OMS 型衰减斜率
\end{itemize}

\paragraph{Gap 9（新增）：缺陷工程的定量 OMS 控制缺失}\label{gap-9ux65b0ux589eux7f3aux9677ux5de5ux7a0bux7684ux5b9aux91cf-oms-ux63a7ux5236ux7f3aux5931}

\begin{itemize}
\tightlist
\item
  \textbf{类型}：缺失连接，\textbf{置信度}：0.80
\item
  \textbf{核心发现}：缺陷类型 \textbf{二分}--缺失配体型（HBDC--额外 OMS--容量 +，R11/\cite{p7}/\cite{p13}/\cite{p192}）vs 溶剂甲酸盐占据型（DMF 分解--占据 OMS--容量 -，Nature Commun 2023）。若忽略类型，任何''缺陷-容量''关联都可能误判（R25/R31）。
\item
  \textbf{验证方案}：HBDC 比例梯度--CO 探针 IR--容量关联；13C NMR 定量甲酸盐缺陷；507 缺陷 MOF 生成工具（\cite{r11s3_e8a9b380a8bb}）预筛
\end{itemize}

\subsubsection{4.2 中等严重程度 Gap（Gap 3-8, 10）}\label{ux4e2dux7b49ux4e25ux91cdux7a0bux5ea6-gapgap-3-8-10}

{\def\LTcaptype{none} % do not increment counter
\begin{longtable}[]{@{}
  >{\raggedright\arraybackslash}p{(\linewidth - 8\tabcolsep) * \real{0.1471}}
  >{\raggedright\arraybackslash}p{(\linewidth - 8\tabcolsep) * \real{0.1765}}
  >{\raggedright\arraybackslash}p{(\linewidth - 8\tabcolsep) * \real{0.1765}}
  >{\raggedright\arraybackslash}p{(\linewidth - 8\tabcolsep) * \real{0.2353}}
  >{\raggedright\arraybackslash}p{(\linewidth - 8\tabcolsep) * \real{0.2647}}@{}}
\toprule\noalign{}
\begin{minipage}[b]{\linewidth}\raggedright
Gap
\end{minipage} & \begin{minipage}[b]{\linewidth}\raggedright
名称
\end{minipage} & \begin{minipage}[b]{\linewidth}\raggedright
类型
\end{minipage} & \begin{minipage}[b]{\linewidth}\raggedright
置信度
\end{minipage} & \begin{minipage}[b]{\linewidth}\raggedright
关键证据
\end{minipage} \\
\midrule\noalign{}
\endhead
\bottomrule\noalign{}
\endlastfoot
Gap 3 & 容量-选择性-再生能 Pareto 前沿未刻画 & 缺失连接 & 0.85 & \cite{p27}（Qst 调控）; \cite{p176}/\cite{p178}（TSA/VPSA CFD） \\
Gap 4 & ML 筛选-实验闭环断裂 + 数据库误差 & 缺失连接 & 0.75 & MP/OQMD 对 MOF 吸附性质 0 覆盖（方法论发现） \\
Gap 5 & DAC（400 ppm）数据稀缺 & 未探索空间 & 0.85 & \cite{p31}/\cite{p50}/\cite{p35}/\cite{p60}/\cite{p54}（综述支持） \\
Gap 6 & OMS 密度-Qst 标度律缺失 & 缺失连接 & 0.72 & \cite{p7}/\cite{p13}/\cite{p192}/\cite{p17}/\cite{p185} \\
Gap 7 & 材料-再生工艺耦合优化缺失 & 缺失连接 & 0.78 & \cite{p176}/\cite{p178}/\cite{p196}/\cite{p35}/\cite{p188} \\
Gap 8 & 杂质气体（NO2/SO2）影响研究稀缺 & 未探索空间 & 0.70 & \cite{p216}（唯一直接证据） \\
Gap 10 & 双金属实际组成 vs 名义比例系统偏差 & 缺失连接 & 0.78 & R22（NMR 非随机配分）+ R23（机械化学可控）+ R24（偏差系统化） \\
\end{longtable}
}

\subsubsection{4.3 优先级排序}\label{ux4f18ux5148ux7ea7ux6392ux5e8f}

{\def\LTcaptype{none} % do not increment counter
\begin{longtable}[]{@{}lllll@{}}
\toprule\noalign{}
排名 & Gap & 严重程度 & 置信度 & 与发现关联 \\
\midrule\noalign{}
\endhead
\bottomrule\noalign{}
\endlastfoot
1 & Gap 1 双金属比例-容量 & 高 & 0.95 & 发现 2（1:1 协同） \\
2 & Gap 2 水-CO2 机理 & 高 & 0.97 & 发现 4（RH 响应） + 辫状链机理 \\
3 & Gap 9 缺陷类型-OMS 定量 & 高 & 0.80 & H3 类型依赖修正 \\
4 & Gap 3 Qst Pareto & 高 & 0.85 & 发现 1（Qst 甜点区） \\
5 & Gap 10 实际组成偏差 & 中--\textgreater 高 & 0.78 & 机械化学/NMR 闭环 \\
6 & Gap 7 材料-工艺耦合 & 中 & 0.78 & 工艺证据充足 \\
7 & Gap 5 DAC & 中 & 0.85 & 综述级证据 \\
8 & Gap 8 杂质气体 & 中 & 0.70 & 单点证据 \\
9 & Gap 6 OMS-Qst 标度律 & 中 & 0.72 & 发现 3 支撑 \\
10 & Gap 4 ML 闭环 & 中 & 0.75 & -- \\
\end{longtable}
}

\begin{center}\rule{0.5\linewidth}{0.5pt}\end{center}

\subsection{五、阶段二：构效关系发现指引}\label{ux4e94ux9636ux6bb5ux4e8cux6784ux6548ux5173ux7cfbux53d1ux73b0ux6307ux5f15}

\subsubsection{5.1 假设生成（5 条）}\label{ux5047ux8bbeux751fux62105-ux6761}

基于 Gap 优先级排序，阶段二生成以下假设：

{\def\LTcaptype{none} % do not increment counter
\begin{longtable}[]{@{}
  >{\raggedright\arraybackslash}p{(\linewidth - 12\tabcolsep) * \real{0.0536}}
  >{\raggedright\arraybackslash}p{(\linewidth - 12\tabcolsep) * \real{0.1071}}
  >{\raggedright\arraybackslash}p{(\linewidth - 12\tabcolsep) * \real{0.1607}}
  >{\raggedright\arraybackslash}p{(\linewidth - 12\tabcolsep) * \real{0.1429}}
  >{\raggedright\arraybackslash}p{(\linewidth - 12\tabcolsep) * \real{0.1607}}
  >{\raggedright\arraybackslash}p{(\linewidth - 12\tabcolsep) * \real{0.1607}}
  >{\raggedright\arraybackslash}p{(\linewidth - 12\tabcolsep) * \real{0.2143}}@{}}
\toprule\noalign{}
\begin{minipage}[b]{\linewidth}\raggedright
\#
\end{minipage} & \begin{minipage}[b]{\linewidth}\raggedright
假设
\end{minipage} & \begin{minipage}[b]{\linewidth}\raggedright
源自 Gap
\end{minipage} & \begin{minipage}[b]{\linewidth}\raggedright
置信度
\end{minipage} & \begin{minipage}[b]{\linewidth}\raggedright
Novelty
\end{minipage} & \begin{minipage}[b]{\linewidth}\raggedright
搜索方式
\end{minipage} & \begin{minipage}[b]{\linewidth}\raggedright
Best Score
\end{minipage} \\
\midrule\noalign{}
\endhead
\bottomrule\noalign{}
\endlastfoot
hypo\_1 & 双金属 MOF-74 金属比例与 CO2 容量呈倒 U 型定量关系 & Gap 1 & 0.88 & 0.82 & bayesian & 0.665 \\
hypo\_2 & 胺功能化 MOF 湿态 CO2 容量随 RH 呈峰值增强（胺型促进 vs OMS 抑制双分支） & Gap 2 & 0.90 & 0.85 & bayesian & \textbf{0.829} \\
hypo\_3 & MOF Qst 在 25--40 kJ/mol 窗口实现容量-选择性-再生能耗 Pareto 最优 & Gap 3 & 0.62 & 0.78 & bayesian & 0.677 \\
hypo\_4 & MOF-74 缺陷类型依赖：缺失配体型 vs 溶剂甲酸盐占据型对容量影响相反 & Gap 9 & \textbf{0.92} & \textbf{0.88} & bayesian & \textbf{0.913} \\
hypo\_5 & MOF 在痕量 NO2/SO2 暴露后 CO2 容量衰减率与 OMS 密度相关 & Gap 8 & 0.63 & \textbf{0.90} & bayesian & 0.734 \\
\end{longtable}
}

合计 165 个候选点探索（\texttt{total\_explored=165}），search\_summary 判定 strong=5。

\subsubsection{5.2 关键发现信号（四类）}\label{ux5173ux952eux53d1ux73b0ux4fe1ux53f7ux56dbux7c7b}

{\def\LTcaptype{none} % do not increment counter
\begin{longtable}[]{@{}
  >{\raggedright\arraybackslash}p{(\linewidth - 2\tabcolsep) * \real{0.5000}}
  >{\raggedright\arraybackslash}p{(\linewidth - 2\tabcolsep) * \real{0.5000}}@{}}
\toprule\noalign{}
\begin{minipage}[b]{\linewidth}\raggedright
类型
\end{minipage} & \begin{minipage}[b]{\linewidth}\raggedright
内容
\end{minipage} \\
\midrule\noalign{}
\endhead
\bottomrule\noalign{}
\endlastfoot
\textbf{正结果} & 双金属倒 U 型（hypo\_1, 0.88）；水 RH 双分支（hypo\_2, 0.90）；缺陷类型二分（hypo\_4, 0.92）--均是置信度 \textgreater= 0.7 且有 p\# 证据链 \\
\textbf{负结果} & MP/OQMD 对 MOF 吸附性质覆盖为 0 --升级为方法论发现（Gap 4 延伸证据） \\
\textbf{异常} & Ni-MOF-74 容量 3.99 vs 8.29 mmol/g 矛盾区间；水-CO2 竞争/协同机理矛盾 \\
\textbf{反例} & 胺功能化 MOF 低湿度容量反升（水促进分支），与''水总是有害''朴素假设相反 \\
\end{longtable}
}

\subsubsection{5.3 定量验证状态}\label{ux5b9aux91cfux9a8cux8bc1ux72b6ux6001}

\begin{itemize}
\tightlist
\item
  \textbf{贝叶斯 vs 随机参照系}：同预算 10 种子公平对比，3:2（bayesian\_wins=3 / random\_wins=2），区分度有限（分数集中 0.54--0.68）。当前发现信号主要来自文献证据结构。
\item
  \textbf{统计显著性}：嵌套 F 检验 hypo\_1 二次 vs 线性（F=9.909, p=0.0254）在 =0.05 下显著。
\item
  \textbf{``候选 vs 经典模型''对比}：因数据不足未达成（达成路径已明确：结构化表格化知识图谱数值 --\textgreater{} 回归核验 --\textgreater{} 经典模型对照）。这是路线 A 复赛需补齐的核心得分点。
\end{itemize}

\begin{center}\rule{0.5\linewidth}{0.5pt}\end{center}

\subsection{六、结论与展望}\label{ux516dux7ed3ux8bbaux4e0eux5c55ux671b}

\subsubsection{6.1 核心结论}\label{ux6838ux5fc3ux7ed3ux8bba}

\begin{enumerate}
\def\labelenumi{\arabic{enumi}.}
\item
  \textbf{MOF CO2 捕获领域知识已高度碎片化但结构化可行}：经 11 轮 546 篇文献的自主体检，Agent 成功构建了 33 条构效关系与 10 项 Research Gap 的知识框架，证明 LLM Agent 在文献驱动科学发现中可扮演''知识组织者 + 矛盾探测器''角色。
\item
  \textbf{三大科学问题构成可直接验证的假说池}：双金属比例-容量定量关系（Gap 1）、水-CO2 机理统一（Gap 2）、缺陷类型-容量方向二分（Gap 9）均已积累充足证据和验证方案，可直接进入实验验证阶段。
\item
  \textbf{方法论发现具有独立价值}：MP/OQMD 对 MOF 吸附性质零覆盖的发现警示了''AI 材料筛选''依赖数据库可能存在的系统性盲区（Gap 4）。
\end{enumerate}

\subsubsection{6.2 局限性}\label{ux5c40ux9650ux6027}

\begin{itemize}
\tightlist
\item
  Gap 2（水-CO2 机理）的辫状链模型仍为 ab initio 预测（2025），缺少独立实验验证
\item
  Gap 1 的倒 U 型规律以 d-d 组合（NiCo/CoMn）为主体证据，s-d 组合（Cu/Mg）单调反例仅 1 个体系
\item
  多数数值（如 Qst、吸附容量）来自文献摘录而非统一条件下的系统测量，存在合成条件/后处理/测量方法的系统性变异性（R26）
\end{itemize}

\subsubsection{6.3 对复赛的展望}\label{ux5bf9ux590dux8d5bux7684ux5c55ux671b}

复赛阶段应聚焦： 1. \textbf{补齐''统计优于前人''的定量验证}（模型对比：候选回归 vs Slack/Vegard 经典模型） 2. \textbf{全量重跑 v2.0 LLM 搜索引导}，产出 LLM 深度参与搜索推理的审计记录 3. \textbf{结构化知识图谱中的量化建模数值表}，为符号回归提供 (x,y) 数据源 4. \textbf{跨主题泛化性验证}：在钙钛矿/热电/正极等 6 个并行主题上检验 Agent 通用性

\begin{center}\rule{0.5\linewidth}{0.5pt}\end{center}

\subsection{证据链}\label{ux8bc1ux636eux94fe}

\begin{quote}
本报告的每个 Gap 编号（Gap 1-10）、每条构效关系（R1-R33）、每个论文引用（p\# / r\# 编号）均可通过以下路径追溯：

\textbf{Gap 编号 / 假设结论 --\textgreater{} \texttt{paper\_summaries.md}（论文完整信息） --\textgreater{} \texttt{search\_log.jsonl}（检索查询 + 数据源） --\textgreater{} \texttt{sciverse\_skill\_log.jsonl}（原始 API 调用记录）}

详细追溯机制见 \texttt{audit\_trail.md}。所有引用均来自真实检索结果或历史记忆文件，零虚构引用。
\end{quote}

\begin{center}\rule{0.5\linewidth}{0.5pt}\end{center}

\emph{本报告由 Pi-Agent（PiAgent）自主生成，过程包含人工审核与格式整理。} \emph{Agent 版本：v2.0 \textbar{} 生成日期：2026-08-02 \textbar{} 文献基础：546 篇（11 轮迭代）}

\section{研究空白（Gap）分析报告 --- MOF 材料用于 CO2 捕获}\label{gap:ux7814ux7a76ux7a7aux767dgapux5206ux6790ux62a5ux544a-mof-ux6750ux6599ux7528ux4e8e-co2-ux6355ux83b7}

\begin{quote}
版本：第二轮更新（2026-08-01） 基础：首轮 152 篇（历史记忆）+ 本轮 Crossref 224 篇 每条 Gap 标注：类型 / 严重程度 / 置信度 / 证据 / 验证方案
\end{quote}

\begin{center}\rule{0.5\linewidth}{0.5pt}\end{center}

\subsection{Gap 1：双金属 MOF 金属比例-容量定量关系缺失}\label{gap:gap-1ux53ccux91d1ux5c5e-mof-ux91d1ux5c5eux6bd4ux4f8b-ux5bb9ux91cfux5b9aux91cfux5173ux7cfbux7f3aux5931}

\begin{itemize}
\tightlist
\item
  \textbf{类型}：缺失连接（定量关系）｜\textbf{严重程度}：高｜\textbf{置信度}：0.95（↑ 自 0.90）
\item
  \textbf{证据}：

  \begin{itemize}
  \tightlist
  \item
    \cite{p65}（10.3390/su17073060）：CoMn-MOF-74 最佳 Co/Mn = \textbf{1:1}，弥补单金属缺陷（废水吸附验证）
  \item
    \cite{p67}（10.9767/bcrec.20658）：Fe/Cu-MOF 双金属 \textbf{\textgreater{} 单金属} Fe-MOF 和 Cu-MOF（1-5 bar 系统评估）
  \item
    \cite{p22}（10.3390/membranes15120385）：Ni-Cu/Co/Zn-MOF-74 系列，Ni-Cu 最优（3188 Barrer / 35.1 选择性）
  \item
    \cite{p24}（10.1016/j.efmat.2023.01.002）：NiCo-MOF-74 微波合成 OMS 增强
  \item
    \cite{p147}（10.1016/j.ces.2016.03.035）：MIL-101(Cr,Mg) 双金属高容量+高选择性
  \item
    历史 \cite{p62}（NiCo 8.30 vs Ni 3.99/Co 5.03）、\cite{p169}（混合金属热力学偏好）
  \end{itemize}
\item
  \textbf{矛盾点}：\cite{p20}（Ni 8.29）与 \cite{p62}（Ni 3.99）数值差异大 → 合成条件敏感；⭐ 九轮新增第三数据点：两阶段离心洗涤优化 Ni-MOF-74 达 \textbf{5.80 mmol/g}（1bar/25°C，显著高于基础合成，Mater 2020）------Ni-MOF-74 容量跨 3.99-8.29 mmol/g 区间，合成/后处理条件主导有效容量
\item
  \textbf{验证方案}：梯度金属比合成（0-100\%，步长 10\%）+ GCMC/DFT 计算，建立''比例→容量/选择性''定量曲线；对倒 U 型峰值位置（1:1 vs 3:1）进行统计检验；\textbf{必须配套 ICP 实测组成 + 后处理标准化（洗涤/激活）以消除 R26 型容量噪声}
\end{itemize}

\subsection{Gap 2：水-CO2 竞争/协同机理矛盾（胺分支十一轮闭环，OMS vs 胺二分机理统一）}\label{gap:gap-2ux6c34-co2-ux7adeux4e89ux534fux540cux673aux7406ux77dbux76feux80faux5206ux652fux5341ux4e00ux8f6eux95edux73afoms-vs-ux80faux4e8cux5206ux673aux7406ux7edfux4e00}

\begin{itemize}
\tightlist
\item
  \textbf{类型}：矛盾结论→机理统一｜\textbf{严重程度}：高｜\textbf{置信度}：0.95 → \textbf{0.97}（十一轮：dobpdc 系多体系实证 + 辫状链原子尺度机理 + IRMOF-74-III 湿态实证）
\item
  \textbf{证据}：

  \begin{itemize}
  \tightlist
  \item
    OMS 分支已解：\cite{p186}（10.1016/j.commatsci.2024.113462）功能化 Mg-MOF-74 不同水蒸气浓度分子模拟；\cite{p182}（10.1016/j.commatsci.2022.111407）配体功能化湿 CO2 吸附；\cite{p196}（10.1016/j.apenergy.2018.09.069）CO2/N2/H2O 突破+循环；历史 \cite{p168}（DOI: 10.1021/acsearthspacechem.7b00142；体系：Al³⁺ 水溶液水解 DFT，机理旁证）（水解离机理）；\cite{r11s4_b2c67eefd8bc}（2016 JMCA：H2O→OH+H 毒化金属中心启动 MOF-74 降解）
  \item
    胺分支新证据：\cite{p210}（10.1016/j.ccst.2024.100356）MOF@MOF 核壳水蒸气亲和调控 → 湿态稳态捕获；\cite{p224}（10.1016/j.seppur.2023.123606）疏水功能化胺浸渍树脂湿空气捕获；\cite{p201}（10.1016/j.jcou.2025.103183）疏水核壳 zeolite@MOF
  \item
    ⭐ 十轮新增（胺分支闭环，待办3）：① 双胺接枝 MOF（diamine-appended）突破实验+协同晶格理论------低 RH 下 CO2 吸附增强、水促进氨基甲酸酯形成（2024, \cite{r10s0_7e6841cd666f}）------\textbf{COF 协同机理（单水吸附位点）成功迁移至胺型 MOF}；② 胺功能化 UiO-66 中\textbf{水桥替代缺陷} boosting CO2 吸附（2021, \cite{r10s0_1fbbe940f717}）------水-缺陷机理交叉；③ Mg-MOF-74+TEPA 后合成功能化干湿双条件性能提升（2017, \cite{r10s0_4eac43c57449}）
  \item
    ⭐⭐ 十一轮新增（R30 普适性检验通过，胺分支机理级闭环）：① \textbf{``water enables efficient CO2 capture''}------氧化抗性 diamine-appended MOF 中水使低分压（≤4 mbar）CO2 捕获成为可能（JACS 2019, \cite{r11s0_5514c3aa655e}，标题即结论）；② 协同吸附剂粗天然气开关式捕获湿态性能增强（Chem Sci 2022, \cite{r11s0_d9371eaa1c08}）；③ \textbf{H2O/CO2 共吸附统计力学+ab initio 预测''辫状链''（braided chain）构型}------H2O 与 CO2 在 Mg2(dobpdc) 孔道交织成链（2025, \cite{r11s2_1fdf67d949bc}）------水促进机理的原子尺度证据；④ IRMOF-74-III-CH2NH2 在 \textbf{65\% RH} 湿态下选择性捕获 CO2（JACS 2014, \cite{r11s4_2c91341a5f47}）------OMS 型之外的湿态实证；⑤ 靶向缺陷对 SP-MOF 湿态烟气 CO2 捕获影响（分子模拟, JCOU 2022, \cite{r11s3_8dfa6887a211}）------缺陷-水交叉独立体系再现
  \item
    ⭐ 九轮新增机理级证据：聚酰亚胺 COF 在 N2/CO2/H2O 突破实验中，\textbf{有限 RH 下 H2O 竞争被协同吸附取代：CO2 容量比干燥高 25\%（343K, 10\% RH）}；机理=CO2 吸附于单一位点吸附水上（FT-IR 证实）；水簇形成后容量损失不可避免；75h 暴露+403K 性能保持（ACS AMI 2023, competitive and cooperative CO2-H2O）
  \item
    历史矛盾：\cite{p16}（水促进，原文待补 DOI）vs \cite{p129}（无影响）vs \cite{p139}（无竞争，原文待补 DOI）vs \cite{p128}（衰减）------\textbf{十一轮结论：矛盾可统一为 OMS 型（水占据 OMS→抑制）vs 胺型（水参与链形成→促进）二分，且有辫状链原子尺度机理支撑}
  \end{itemize}
\item
  \textbf{验证方案}：统一 0-90\% RH breakthrough 矩阵（OMS 型 vs 胺型对照），量化''胺型峰值 RH 位置''与''OMS 型衰减斜率''；将 COF 协同机理（单水吸附位点）迁移检验至胺功能化 MOF；\textbf{十一轮新增可用材料库：mmen/dmpn/2-ampd/dap-Mg2(dobpdc) 系列（\cite{r11s1_a7d3df468203}, \cite{r11s0_9932dacb5791}, \cite{r11s1_ff6147b117a5}）+ IRMOF-74-III-CH2NH2（\cite{r11s4_2c91341a5f47}）}
\end{itemize}

\subsection{Gap 3：容量-选择性-再生能 Pareto 前沿未刻画}\label{gap:gap-3ux5bb9ux91cf-ux9009ux62e9ux6027-ux518dux751fux80fd-pareto-ux524dux6cbfux672aux523bux753b}

\begin{itemize}
\tightlist
\item
  \textbf{类型}：缺失连接｜\textbf{严重程度}：高｜\textbf{置信度}：0.85（↑ 自 0.80）
\item
  \textbf{证据}：\cite{p27}（10.1016/j.jcis.2021.12.163）MOF-74(Ni) Qst 调控；\cite{p117}（10.1039/d0dt01784a）Qst 计算指南；\cite{p129}（10.3390/jcs5040102）胺化 SBA-15 Qst；\cite{p116} 锂掺杂 MIL-101 异常 Qst；历史 \cite{p27}（Qst 29 甜点）、\cite{p118}（25--40 kJ/mol 窗口；同 \cite{p27}，DOI: 10.1016/j.jcis.2021.12.163）
\item
  \textbf{验证方案}：50+ 材料三维性能图（容量×选择性×再生能），识别 Pareto 前沿；用 Qst 甜点区假设（29-40 kJ/mol）做约束优化
\end{itemize}

\subsection{Gap 4：ML 筛选-实验闭环断裂 + 数据库误差}\label{gap:gap-4ml-ux7b5bux9009-ux5b9eux9a8cux95edux73afux65adux88c2-ux6570ux636eux5e93ux8befux5dee}

\begin{itemize}
\tightlist
\item
  \textbf{类型}：缺失连接｜\textbf{严重程度}：中｜\textbf{置信度}：0.75
\item
  \textbf{证据}：\cite{p191}（10.3390/chemengineering9040080）沸石 ML 通用预测（5700+ 数据点）；历史 \cite{p17}（LitMOF 数据库误差）；MP/OQMD 对 MOF 吸附性质 0 覆盖（首轮发现）
\item
  \textbf{验证方案}：构建 MOF 专属吸附数据集，训练可迁移 ML 模型，与 GCMC 交叉验证
\end{itemize}

\subsection{Gap 5：DAC（400 ppm）数据稀缺}\label{gap:gap-5dac400-ppmux6570ux636eux7a00ux7f3a}

\begin{itemize}
\tightlist
\item
  \textbf{类型}：未探索空间｜\textbf{严重程度}：中｜\textbf{置信度}：0.85（↑ 自 0.80）
\item
  \textbf{证据}：\cite{p31}（10.1016/j.envres.2024.119985）MOF DAC 综述；\cite{p50}（10.3390/molecules30143048）DAC 纳米材料综述（强调水敏感性）；\cite{p35}（10.1016/j.cej.2025.170309）MOF-碳纤维电驱动再生 DAC；\cite{p60}（TVSA 高纯 CO2）
\item
  \textbf{验证方案}：400 ppm 与 15\% 浓度下的容量/动力学对比实验；MOF DAC 循环稳定性长测
\end{itemize}

\subsection{Gap 6：OMS 密度-Qst 标度律缺失}\label{gap:gap-6oms-ux5bc6ux5ea6-qst-ux6807ux5ea6ux5f8bux7f3aux5931}

\begin{itemize}
\tightlist
\item
  \textbf{类型}：缺失连接｜\textbf{严重程度}：中｜\textbf{置信度}：0.72（↑ 自 0.70）
\item
  \textbf{证据}：\cite{p7}/\cite{p13}（缺陷→额外 OMS）、\cite{p192}（10.1016/j.cej.2023.144052）缺陷丰富 Mg-MOF-74 增强吸附；\cite{p17}（10.1021/acs.langmuir.5c04277）M-MOF-74 结合势能景观；\cite{p185} M-MOF-74 分子模拟
\item
  \textbf{验证方案}：用缺陷浓度梯度（HBDC/DHBDC 比例）定量 OMS 密度，测量 Qst 标度律斜率
\end{itemize}

\subsection{Gap 7：材料-再生工艺耦合优化缺失}\label{gap:gap-7ux6750ux6599-ux518dux751fux5de5ux827aux8026ux5408ux4f18ux5316ux7f3aux5931}

\begin{itemize}
\tightlist
\item
  \textbf{类型}：缺失连接｜\textbf{严重程度}：中｜\textbf{置信度}：0.78（↑ 自 0.70）
\item
  \textbf{证据}：\cite{p176}（10.1016/j.enconman.2017.11.010）Mg-MOF-74 TSA 工艺；\cite{p178}（10.1016/j.apenergy.2017.10.098）Mg-MOF-74 VPSA CFD；\cite{p196} 突破+循环；\cite{p35} 电驱动再生；\cite{p188} TSA 动态
\item
  \textbf{验证方案}：材料参数（Qst、容量、动力学）× 工艺参数（温度、真空度、循环时间）联合优化 DOE
\end{itemize}

\subsection{Gap 8（新增）：杂质气体（NO2/SO2）对 MOF CO2 捕获的影响研究稀缺}\label{gap:gap-8ux65b0ux589eux6742ux8d28ux6c14ux4f53no2so2ux5bf9-mof-co2-ux6355ux83b7ux7684ux5f71ux54cdux7814ux7a76ux7a00ux7f3a}

\begin{itemize}
\tightlist
\item
  \textbf{类型}：未探索空间｜\textbf{严重程度}：中｜\textbf{置信度}：0.70
\item
  \textbf{证据}：\cite{p216}（10.3390/gases6020024）ELM-11 暴露 1000 ppm NO2 后 CO2 吸附衰减（唯一直接证据）；\cite{p52}/\cite{p55}（低纯度 CO2 压缩能耗）
\item
  \textbf{验证方案}：对 MOF-74/ZIF/MIL 系列做 NO2/SO2 暴露矩阵（浓度×时间），测量容量衰减率与结构变化（PXRD/FTIR）
\end{itemize}

\subsection{Gap 9（新增）：缺陷工程的定量 OMS 控制缺失}\label{gap:gap-9ux65b0ux589eux7f3aux9677ux5de5ux7a0bux7684ux5b9aux91cf-oms-ux63a7ux5236ux7f3aux5931}

\begin{itemize}
\tightlist
\item
  \textbf{类型}：缺失连接｜\textbf{严重程度}：中 → \textbf{高}｜\textbf{置信度}：0.65 → 0.75 → 0.78 → \textbf{0.80}（十一轮：缺陷-水交叉旁证 + 507 缺陷 MOF 高通量工具）
\item
  \textbf{证据}：

  \begin{itemize}
  \tightlist
  \item
    \cite{p192} 缺陷丰富 vs \cite{p13} 定向缺陷（HBDC 比例），但均未给出 OMS 密度定量值
  \item
    ⭐ 九轮新增（缺陷类型二分）：MOF-74 存在\textbf{溶剂衍生甲酸盐缺陷}（DMF 分解的甲酸配位金属）------甲酸盐额外 M-O 键\textbf{部分消除 OMS，N2/CO2 吸附随缺陷浓度定量下降}（Nature Commun 2023, solvent-derived defects suppress adsorption in MOF-74；多维固态 NMR + DFT 定位甲酸盐构型）
  \item
    ⭐ 十轮新增（类型依赖旁证）：④ Ru-MOF 不同缺陷类型（5-X-ipH2, X=OH/H/NH2/Br）对孔隙率与吸附影响不同（2016, \cite{r10s1_821919f491c8}）；⑤ 缺陷工程 Ni-MOF-74 配体功能化------CO2 吸附取决于片段数量与功能基团类型/位置/大小（2022, \cite{r10s2_0341d732c860}）；⑥ 胺功能化 UiO-66 水桥替代缺陷增强吸附（2021, \cite{r10s0_1fbbe940f717}）------缺陷-水交叉机理------\textbf{与 R1/R11（缺失配体→容量↑）形成矛盾，必须区分缺陷类型}
  \item
    ⭐ 十一轮新增（交叉+工具）：⑦ 靶向缺陷对 SP-MOF 湿态烟气 CO2 捕获的影响（分子模拟, JCOU 2022, \cite{r11s3_8dfa6887a211}）------缺陷-湿度交叉在独立体系再现（R30 普适性）；⑧ \textbf{507 个含缺失配体点缺陷 MOF 高通量生成工具 + 最大缺陷浓度概念}（JPCL 2023, \cite{r11s3_e8a9b380a8bb}）------缺陷类型-浓度模拟数据库为 Gap 9 定量验证提供计算基础设施
  \item
    ⭐ 九轮新增（定量方法）：UiO 系列缺失配体缺陷定量新流程（M-200 vs M-400 激活对比 + 封端剂量化），在 UiO-66/67/68/69 验证有效（Chem Mater 2023, quantification of linker defects in UiO-type MOFs）
  \end{itemize}
\item
  \textbf{严重程度升级理由}：缺陷类型（缺失配体 vs 溶剂占据）对 CO2 吸附方向相反------若忽略类型，任何''缺陷-容量''关联都可能误判；H3 假设（缺陷密度-容量线性正相关）必须按缺陷类型修正
\item
  \textbf{验证方案}：HBDC 比例梯度（0-50\%）→ 定量 OMS 密度（CO 探针 IR）→ 容量/Qst 关联；\textbf{并行追踪溶剂衍生甲酸盐缺陷}（13C NMR 定量，参考 Nature Commun 2023 方法）；区分''缺失配体型''与''溶剂占据型''两种缺陷通道；\textbf{十一轮新增：可复用 507 缺陷 MOF 生成工具（\cite{r11s3_e8a9b380a8bb}）预筛缺陷类型-浓度组合，再实验验证}
\end{itemize}

\begin{center}\rule{0.5\linewidth}{0.5pt}\end{center}

\subsection{优先级排序（下一轮执行）}\label{gap:ux4f18ux5148ux7ea7ux6392ux5e8fux4e0bux4e00ux8f6eux6267ux884c}

{\def\LTcaptype{none} % do not increment counter
\begin{longtable}[]{@{}
  >{\raggedright\arraybackslash}p{(\linewidth - 8\tabcolsep) * \real{0.1395}}
  >{\raggedright\arraybackslash}p{(\linewidth - 8\tabcolsep) * \real{0.1163}}
  >{\raggedright\arraybackslash}p{(\linewidth - 8\tabcolsep) * \real{0.2093}}
  >{\raggedright\arraybackslash}p{(\linewidth - 8\tabcolsep) * \real{0.1860}}
  >{\raggedright\arraybackslash}p{(\linewidth - 8\tabcolsep) * \real{0.3488}}@{}}
\toprule\noalign{}
\begin{minipage}[b]{\linewidth}\raggedright
排名
\end{minipage} & \begin{minipage}[b]{\linewidth}\raggedright
Gap
\end{minipage} & \begin{minipage}[b]{\linewidth}\raggedright
严重程度
\end{minipage} & \begin{minipage}[b]{\linewidth}\raggedright
置信度
\end{minipage} & \begin{minipage}[b]{\linewidth}\raggedright
与历史发现关联
\end{minipage} \\
\midrule\noalign{}
\endhead
\bottomrule\noalign{}
\endlastfoot
1 & Gap 1 双金属比例-容量 & 高 & 0.95 & 发现 2（1:1 协同）直接支撑 \\
2 & Gap 2 水-CO2 机理 & 高 & 0.97（↑0.95） & 发现 4（RH 响应）直接支撑；COF 协同机理旁证（九轮） \\
3 & Gap 9 缺陷类型-OMS 定量 & 高 & 0.80（↑0.78） & H3 直接相关（十轮修正为类型依赖）；三方独立证据 \\
4 & Gap 3 Qst Pareto & 高 & 0.85 & 发现 1（Qst 甜点区）直接支撑 \\
5 & Gap 10 实际组成偏差 & 中→高 & 0.78 & 机械化学/NMR 证据闭环（六轮） \\
6 & Gap 7 材料-工艺耦合 & 中 & 0.78 & 工艺维度新证据充足 \\
7 & Gap 5 DAC & 中 & 0.85 & 综述级证据充分 \\
8 & Gap 8 杂质气体 & 中 & 0.70 & 新维度，单点证据 \\
9 & Gap 6 OMS-Qst 标度律 & 中 & 0.72 & 发现 3 支撑 \\
10 & Gap 4 ML 闭环 & 中 & 0.75 & --- \\
\end{longtable}
}

\begin{center}\rule{0.5\linewidth}{0.5pt}\end{center}

\subsection{第四轮更新（2026-08-01，双金属协同专题 +84 篇）}\label{gap:ux7b2cux56dbux8f6eux66f4ux65b02026-08-01ux53ccux91d1ux5c5eux534fux540cux4e13ux9898-84-ux7bc7}

\subsubsection{Gap 1 更新：双金属比例-性能曲线''倒 U 型非普适性''研究空白}\label{gap:gap-1-ux66f4ux65b0ux53ccux91d1ux5c5eux6bd4ux4f8b-ux6027ux80fdux66f2ux7ebfux5012-u-ux578bux975eux666eux9002ux6027ux7814ux7a76ux7a7aux767d}

\begin{itemize}
\tightlist
\item
  \textbf{置信度}：0.90 → \textbf{0.93}（新增 8 项直接证据）
\item
  \textbf{新增证据}：

  \begin{itemize}
  \tightlist
  \item
    scv:073a40bcc5cb：Ni₀.₃₇Co₀.₆₃-MOF-74 中间比例 \textgreater{} 所有单金属（x=0.37 实测）
  \item
    scv:66fc8c917582：Ni(x)M(1-x)-ITHDs 临界掺杂量依赖组合（NiZn\textgreater0.2 vs NiCo≈0.1）
  \item
    scv:ecfe34f94197：\textbf{Cu/Mg-MOF-74 吸附容量随 Mg 比例单调增（非倒 U），光催化倒 U 峰值 0.6/0.4}------最优比例性质依赖
  \item
    scv:1a24bed1a447：dobpdc 双金属组成受离子半径控制（实际比例可控）
  \item
    scv:39ca0a621b3c / scv:309c6ba090d4：金属离子浸渍/碱金属掺杂增强 Qst+选择性+容量
  \end{itemize}
\item
  \textbf{细化结论}：倒 U 型适用于 d-d 组合（NiCo/CoMn）的吸附容量；但不普适（Cu/Mg 吸附单调）；最优比例随目标性质（吸附 vs 催化）漂移
\item
  \textbf{第五轮新证据（轨道类型假设）}：Cu/Mg-MOF-74（s-d 组合）吸附单调增再次验证（scv:ecfe34f94197 重检索命中，9.21 mmol/g@0.1/0.9）；Mg-MOF-74 部分取代 Co/Ni（s-d 组合）------轨道类型（s vs d）决定曲线形状假设获支撑
\item
  \textbf{置信度}：0.93 → \textbf{0.94}
\item
  \textbf{验证方案升级}：梯度合成需覆盖 (a) 多金属组合（d-d、s-d、s-p 各至少 2 组）(b) 多性质维度（容量、选择性、Qst、催化活性）(c) 比例步长 5-10\%，以绘制''金属组合 × 性质 × 比例''三维响应面
\item
  \textbf{第六轮新增（组成可控性解锁）}：机械化学球磨法以固体配位前驱体实现预定 1:1 化学计量，12 种双金属 MOF-74（含 MgZn/MgCo/MgNi/MgCu 全部 s-d 组合）------轨道类型-曲线形状假说（R19）的实验验证路径从''不可行''变为''可执行''
\item
  \textbf{置信度}：0.94 → \textbf{0.95}
\end{itemize}

\subsection{Gap 10（新增）：双金属 MOF 实际组成 vs 名义比例偏差的系统量化缺失}\label{gap:gap-10ux65b0ux589eux53ccux91d1ux5c5e-mof-ux5b9eux9645ux7ec4ux6210-vs-ux540dux4e49ux6bd4ux4f8bux504fux5deeux7684ux7cfbux7edfux91cfux5316ux7f3aux5931}

\begin{itemize}
\tightlist
\item
  \textbf{类型}：缺失连接｜\textbf{严重程度}：中→高｜\textbf{置信度}：0.60 → 0.70 → \textbf{0.78}（第六轮新增 2 项决定性证据）
\item
  \textbf{证据}：scv:1a24bed1a447（dobpdc 组成受离子半径控制，传统盐直接反应无法得到 1:1）；多数研究只报投料比不报骨架实测比（ICP-MS 缺失）；第五轮新增：Mg-MOF-74 部分取代 Co/Ni 一锅法中''反应温度比溶剂组成更能影响最终金属组成''；MgCu-MOF-74 一锅法采用 ICP-AES 实测组成（验证方法示范）；\textbf{第六轮新增}：① Mg₁₋ₓNiₓ-MOF-74 固态 NMR 配分解析揭示\textbf{全部 8 种 Mg/Ni 原子级排列、驳斥溶液合成随机配分假设}（偏差有系统结构根源，非随机噪声）；② 机械化学球磨法以固体配位前驱体实现\textbf{预定 1:1 化学计量}（12 种双金属 MOF-74，含全部 s-d 组合）
\item
  \textbf{严重程度升级理由}：配分非随机性意味着名义-实际偏差是\textbf{系统性}的（由合成动力学决定），若不经校正，所有已报道的比例-性能曲线（包括倒 U 型证据）都可能存在系统偏移 → 影响 Gap 1 结论的定量可靠性
\item
  \textbf{验证方案}：对已报道双金属 MOF 做 ICP-MS + XPS 组成再分析，构建''名义比例→实际比例''映射库（含温度/溶剂合成参数维度），校正所有比例-性能曲线；\textbf{升级路径}：采用机械化学前驱体法直接合成预定化学计量的双金属 MOF-74（消除偏差），用固态 NMR 验证配分模式，建立''干净''的比例-性能曲线基准
\end{itemize}

\printbibliography[title=参考文献]
\end{document}
